\subsection{Cấu trúc đường ống 5 tầng co giãn của \texttt{posit\_add}}

Bộ cộng dấu phẩy động Posit (\texttt{posit\_add}) được tổ chức thành đường ống
5 tầng co giãn \texttt{S1..S5} nhằm cân bằng độ sâu logic ở mọi tầng và đáp ứng
mục tiêu tần số từ $100\text{ MHz}$ trở lên trên Artix-7: mỗi tầng chỉ chứa tối
đa một phép toán số học hoặc một bộ dịch lớn, giúp năm tầng \texttt{S1..S5}
đạt độ trễ tổ hợp sau tổng hợp nằm trong khoảng $2{,}958\text{--}4{,}851\text{ ns}$
(Bảng~\ref{tab:f2f-stages}) và đóng định thời định tuyến trên MAC-F ở
$F_{\max,\mathrm{R2R}}=154{,}42\text{ MHz}$ (Bảng~\ref{tab:fpga-synth-vs-route}).
Toàn bộ mạch tuần tự và tín hiệu bắt tay \texttt{valid/ready} được tập trung tại
mô-đun đỉnh \texttt{posit\_add.sv}, bao bọc năm khối tổ hợp thuần túy từ S1 đến S5
(Hình~\ref{fig:adder-stages}).

Ở cơ chế co giãn, tín hiệu \texttt{ready} truyền ngược từ S5 về S1: một slot rỗng
có thể nhận ngay, còn slot đang giữ dữ liệu hợp lệ chỉ tiến lên khi tầng kế tiếp
sẵn sàng; nếu bị chặn (\texttt{out\_ready=0}), toàn bộ dữ liệu số học và
metadata (dấu, cờ \texttt{Zero/NaR}, phần dư và \texttt{flags}) được giữ đồng bộ
với nhau vì chỉ lệch một trường cũng làm sai lệch phép đóng gói cuối.

\rtlfig{Datapath mô-đun \sig{posit\_add}. Cột trái là xử lý tổ hợp; cột phải liệt kê các trường chính được chốt sau mỗi tầng.}{fig:adder-stages}{
\node[pin] (in) at (-2.3,.28) {Hai số hạng: frac 26, $sf$ 11, $s,Z,N$ mỗi số};
\node[head] at (4.65,0) {Ngân hàng thanh ghi};
\node[unit,text width=73mm] (s1) at (-2.3,-1.42) {\hw{S1: so và đổi vai}\\
 ưu tiên số khác 0; so $sf$, rồi so frac\\
 hai hiệu scale signed 12 bit $\to |\Delta sf|$; kẹp nếu $\Delta>28$};
\node[ff,text width=50mm] (r1) at (4.65,-1.42) {\hw{S1}\\frac X/Y: 26+26; shift: 5\\kẹp, subtract, $sf_{\mathrm{base}}$, dấu/cờ};
\draw[wire] (in)--(s1);
\draw[wire] (s1)--(r1);
\node[unit,text width=73mm] (s2) at (-2.3,-3.96) {\hw{S2: căn Y vào cửa sổ 31 bit}\\
 $\{0,1,f_Y,3'b0\}\gg\mathrm{shift}$\\
 mặt nạ đuôi $\to$ jam ở \sig{Y[0]} và \sig{alignment\_tail}};
\node[ff,text width=50mm] (r2) at (4.65,-3.96) {\hw{S2}\\frac X: 26; Y: 31; tail: 1\\subtract, $sf_{\mathrm{base}}$, dấu/cờ};
\draw[wire] (r1.south)--++(0,-.54)-|(s2.north);
\draw[wire] (s2)--(r2);
\node[unit,text width=73mm] (s3) at (-2.3,-6.51) {\hw{S3: cộng/trừ độ lớn}\\
 $X=\{0,1,f_X,3'b0\}$; $M=X\pm Y$\\
 song song tính \sig{sf\_inc} $=sf_{\mathrm{base}}+1$};
\node[ff,text width=50mm] (r3) at (4.65,-6.51) {\hw{S3}\\magnitude: 31; tail: 1\\\sig{sf\_inc}: 11, dấu/cờ};
\draw[wire] (r2.south)--++(0,-.54)-|(s3.north);
\draw[wire] (s3)--(r3);
\node[unit,text width=73mm] (s4) at (-2.3,-9.06) {\hw{S4: LOD ưu tiên và seed}\\
 $M[30:0]\to$ \sig{leading\_count[4:0]}, \sig{nonzero}\\
 song song ghép seed 30 bit và \sig{sticky\_raw}};
\node[ff,text width=50mm] (r4) at (4.65,-9.06) {\hw{S4}\\seed: 30; count: 5; nonzero\\sticky, \sig{sf\_inc}, dấu/cờ};
\draw[wire] (r3.south)--++(0,-.54)-|(s4.north);
\draw[wire] (s4)--(r4);
\node[unit,text width=73mm] (s5) at (-2.3,-11.61) {\hw{S5: chuẩn hóa và phân loại}\\
 dịch trái seed theo count; lấy 29 bit cao\\
 $sf=\sig{sf\_inc}-\mathrm{count}$; xác định Zero/NaR};
\node[ff,text width=50mm] (r5) at (4.65,-11.61) {\hw{S5}\\\sig{frac\_hi[28:0]}, $sf$, sticky\\dấu, Zero/NaR và flags};
\draw[wire] (r4.south)--++(0,-.54)-|(s5.north);
\draw[wire] (s5)--(r5);
\node[pin] (o) at (4.65,-13.08) {Tới Packer: \sig{frac = \{frac\_hi,24'b0\}}};
\draw[wire] (r5)--(o);
}



Toàn bộ luồng tính toán năm tầng của \texttt{posit\_add} được tóm tắt
trong Algorithm~\ref{alg:posit-add}, đối chiếu trực tiếp với các khối tổ
hợp từ S1 đến S5 trong RTL.

\begin{reportpseudocode}{alg:posit-add}{Cộng Posit(32,3) đường ống 5 tầng co giãn (Five-Stage Elastic Addition)}
  {Hai số hạng chuẩn hóa $A,B$ (với $s, sf[10:0], f[25:0], z, \eta$) và cờ $\mathrm{flags_{in}}[4:0]$.}
  {Kết quả chuẩn hóa $(s_{\mathrm{out}}, sf[10:0], f_{\mathrm{out}}[52:0], \mathrm{sticky}, z_{\mathrm{out}}, \eta, \mathrm{flags_{out}}[4:0])$.}
  \begin{tabularx}{\linewidth}{@{}r >{\raggedright\arraybackslash}X >{\raggedright\arraybackslash}p{.33\linewidth}@{}}
    \algline{1}{\textbf{S1} \quad $d_{AB}\gets sf_A-sf_B$, $d_{BA}\gets sf_B-sf_A$}{Tính hai hiệu scale mở rộng dấu 12 bit song song.}
    \algline{2}{$\mathrm{choose\_a}\gets \neg z_A\land \bigl(z_B\lor (d_{AB}\ge 0 \land (sf_A\ne sf_B \lor f_A\ge f_B))\bigr)$}{Chọn toán hạng có độ lớn lớn hơn (ưu tiên khác 0).}
    \algline{3}{$(X,Y)\gets (A,B)$ \textbf{if} $\mathrm{choose\_a}$ \textbf{else} $(B,A)$, $\delta \gets \min(|\Delta sf|,29)$}{Hoán đổi toán hạng; xác định độ dịch $\delta$ và cờ kẹp $|\Delta sf|>28$.}
    \algline{4}{$\mathrm{sub}\gets \neg(z_A\lor z_B)\land (s_A\ne s_B)$, $\eta\gets \eta_A\lor\eta_B\lor\mathrm{flags_{in}}[4]$}{Xác định phép trừ hiệu dụng và lan truyền cờ NaR.}
    \algline{5}{\textbf{S2} \quad $w_Y\gets \{1'b0, 1'b1, f_Y[25:0], 3'b0\}$ \textbf{if} $\neg(z_A\lor z_B)$ \textbf{else} $0$}{Mở rộng toán hạng nhỏ thành từ 31 bit có bit ẩn $1$.}
    \algline{6}{$y_{\mathrm{align}}[30:1]\gets (w_Y\gg \delta)[30:1]$, $\mathrm{tail}\gets \mathrm{clamped}\lor \mathrm{raw\_tail}$}{Dịch phải 31 bit và gom mặt nạ nhiệt kế cho bit rơi.}
    \algline{7}{$y_{\mathrm{align}}[0]\gets ((w_Y\gg \delta)[0]\lor \mathrm{tail})$ \textbf{if} $\neg(z_A\lor z_B)$ \textbf{else} $0$}{Chèn bit jam tại LSB khi dịch rơi bit hoặc kẹp $\delta$.}
    \algline{8}{\textbf{S3} \quad $w_X\gets \{1'b0, 1'b1, f_X[25:0], 3'b0\}$ \textbf{if} $\neg(z_A\land z_B)$ \textbf{else} $0$}{Đưa toán hạng lớn về dạng 31 bit để cộng/trừ.}
    \algline{9}{$M\gets w_X-y_{\mathrm{align}}$ \textbf{if} $\mathrm{sub}$ \textbf{else} $w_X+y_{\mathrm{align}}$, $sf_{\mathrm{inc}}\gets sf_X+1$}{Cộng/trừ 31 bit và cộng bù trước $1$ cho scale.}
    \algline{10}{\textbf{S4} \quad $(K,\mathrm{nz})\gets \mathrm{LOD}(M[30:0])$, $\mathrm{lost}\gets M[30]\land M[0]$}{Dò bit 1 dẫn đầu thống nhất 31 bit; kiểm tra carry bị mất.}
    \algline{11}{$\mathrm{seed}[29:0]\gets \{M[29:2], M[1]\lor\mathrm{lost}, M[0]\}$, $\mathrm{sticky}_{\mathrm{raw}}\gets \mathrm{tail}\lor\mathrm{lost}$}{Tạo hạt dịch 30 bit và cập nhật sticky thô.}
    \algline{12}{\textbf{S5} \quad $f_{\mathrm{hi}}[28:0]\gets (\mathrm{seed}\ll K)[29:1]$, $sf\gets sf_{\mathrm{inc}}-K$}{Dịch trái 30 bit chuẩn hóa và cập nhật scale đồng nhất.}
    \algline{13}{$z_{\mathrm{out}}\gets \neg\eta\land ((z_A\land z_B)\lor \neg\mathrm{nz})$, $f_{\mathrm{out}}\gets \{f_{\mathrm{hi}}, 24'b0\}$}{Phân loại Zero/NaR và nối 24 bit 0 tạo ngõ ra 53 bit.}
    \algline{14}{\textbf{return} $s_{\mathrm{out}}, sf, f_{\mathrm{out}}, \mathrm{sticky}_{\mathrm{raw}}, z_{\mathrm{out}}, \eta, \mathrm{flags}_{\mathrm{out}}$}{Chốt thanh ghi S5 kèm cập nhật bit \texttt{inexact}.}
  \end{tabularx}
\end{reportpseudocode}

\subsection{Tầng S1: Chuẩn bị, so sánh độ lớn và hoán đổi toán hạng}

Ở đầu vào MAC-F, $X$ và $Y$ của bộ cộng là hai số hạng được sắp lại theo độ lớn
($|X|\ge|Y|$); ký hiệu này dùng trong phạm vi bộ cộng và khác với vai $X/Y$ của
khối OPS ở bộ nhân. Mỗi số hạng đưa vào S1 gồm fraction 26 bit, scale signed 11 bit,
dấu và cờ đặc biệt. Tầng S1 (\texttt{posit\_add\_prepare\_comb}, Hình~\ref{fig:adder-s1})
tính song song hai hiệu scale mở rộng dấu 12 bit (\texttt{diff\_ab}, \texttt{diff\_ba}),
so sánh ngang bằng scale (\texttt{sf\_equal}) và so sánh fraction (\texttt{frac\_a\_ge\_b})
để tạo tín hiệu điều khiển \texttt{choose\_a\_cond} ưu tiên toán hạng khác 0.

Cờ kẹp \texttt{clamp\_ab/\!ba} kiểm tra đồng thời nhóm bit cao \texttt{diff[10:5]}
khác 0 hoặc \texttt{diff[4:0] > 28} để nhận biết $|\Delta sf|>28$ mà không bị
cuộn modulo 5 bit khi chênh lệch scale lớn. Để giảm fanout và tránh nghẽn định
tuyến trên đường tráo toán hạng $62$ nhánh, RTL tách tín hiệu chọn thành bốn lưới
có thuộc tính \texttt{dont\_touch}: \texttt{choose\_a\_frac0..2} điều khiển riêng
ba lát fraction \texttt{[8:0]}, \texttt{[17:9]}, \texttt{[25:18]} và
\texttt{choose\_a\_meta} điều khiển \texttt{sign\_base}/\texttt{sf\_base}, bảo toàn
nguyên vẹn ranh giới thanh ghi và số chu kỳ đường ống.

\rtlfig{Datapath chi tiết tầng S1 (\sig{posit\_add\_prepare\_comb}): hai bộ trừ scale 12 bit song song, phát hiện kẹp $|\Delta sf|>28$ và bốn nhánh \sig{choose\_a\_*} cân bằng tải.}{fig:adder-s1}{
\node[pin] (ina) at (-5.05,.82) {\sig{a\_sf},\,\sig{b\_sf[10:0]}};
\node[pin] (inb) at (.20,.82) {\sig{a\_frac},\,\sig{b\_frac[25:0]}};
\node[pin] (inc) at (5.35,.82) {dấu, cờ \sig{Z/N}, \sig{flags\_in[4:0]}};

\node[unit,text width=40mm] (sub) at (-5.05,-.92) {\hw{Trừ scale 12 bit song song}\\
 $\sig{diff\_ab}=\{a_{10},a_{sf}\}-\{b_{10},b_{sf}\}$\\
 $\sig{diff\_ba}=\{b_{10},b_{sf}\}-\{a_{10},a_{sf}\}$\\
 $\sig{sf\_a\_lt\_b}=\sig{diff\_ab}[11]$};
\node[unit,text width=45mm] (cmp) at (.20,-.92) {\hw{So sánh độ lớn toán hạng}\\
 $\sig{sf\_equal}=(a_{sf}==b_{sf})$\\
 $\sig{frac\_a\_ge\_b}=(a_{\mathrm{frac}}\ge b_{\mathrm{frac}})$};
\node[unit,text width=38mm] (exc) at (5.35,-.92) {\hw{Cờ ngoại lệ \& phép tính}\\
 $\sig{both\_nz}=\neg a_Z\land\neg b_Z$\\
 $\sig{sub}=\sig{both\_nz}\land(a_s\ne b_s)$\\
 \sig{both\_zero}, \sig{is\_nar}, \sig{flags\_lo}};

\draw[wire] (ina)--(sub);
\draw[wire] (inb)--(cmp);
\draw[wire] (inc)--(exc);

\node[unit,text width=40mm] (clp) at (-5.05,-3.92) {\hw{Phát hiện kẹp \& dịch}\\
 $\sig{clamp\_ab/\!ba}: |\Delta sf|>28$\\
 chọn theo \sig{sf\_a\_lt\_b}:\\
 \sig{shift\_amount[4:0]}, \sig{clamped}};
\node[pick,text width=47mm] (sel) at (.20,-3.92) {\hw{Chọn $X{\ge}Y$ \& tách 4 tải}\\
 $\sig{choose\_a\_cond}=\neg a_Z\land(b_Z\lor\dots)$\\
 4 lưới \sig{dont\_touch}:\\
 \sig{choose\_a\_frac0..2}, \sig{meta}};

\draw[wire] (sub)--(clp);
\draw[wire] (cmp)--(sel);

\coordinate (gx) at (-1.50,-2.65);
\node[tapdot] at (-5.05,-2.65) {};
\draw[meta] (-5.05,-2.65) -- (gx) -- (sel.north -| gx);
\node[tag] at (-3.15,-2.65) {\sig{sf\_a\_lt\_b}};

\coordinate (gz) at (1.90,-2.65);
\node[tapdot] at (5.35,-2.65) {};
\draw[meta] (5.35,-2.65) -- (gz) -- (sel.north -| gz);
\node[tag] at (3.65,-2.65) {$a_Z, b_Z$};

\node[unit,text width=51mm] (muxxy) at (.20,-6.22) {\hw{MUX tráo 3 lát \& metadata}\\
 $[8:0]\gets\sig{frac0}$, $[17:9]\gets\sig{frac1}$, $[25:18]\gets\sig{frac2}$\\
 xuất \sig{x\_frac}, \sig{y\_frac}, \sig{sign\_base}, \sig{sf\_base}};
\draw[meta] (sel)--(muxxy);
\node[tag,anchor=west] at (.38,-5.10) {4 nhánh \sig{dont\_touch}};

\node[ff,text width=140mm] (r1) at (.15,-7.95) {\hw{S1}\enspace
 \sig{s1\_x\_frac\_q}, \sig{s1\_y\_frac\_q[25:0]} \enspace|\enspace
 \sig{s1\_shift\_q[4:0]}, \sig{s1\_clamped\_q}\\
 \sig{s1\_sub\_q}, \sig{s1\_sign\_q}, \sig{s1\_sf\_q[10:0]} \enspace|\enspace
 \sig{s1\_both\_nz\_q}, \sig{s1\_both\_zero\_q}, \sig{s1\_nar\_q}, \sig{s1\_flags\_q[3:0]}};

\draw[wire] (clp.south)--(r1.north -| clp);
\draw[wire] (muxxy.south)--(r1.north -| muxxy);
\draw[wire] (exc.south)--(r1.north -| exc);

\node[pin] (o) at (.15,-9.12) {Tới S2 (\sig{posit\_add\_align\_comb}): căn cửa sổ 31 bit, tạo \sig{jam} và \sig{alignment\_tail}};
\draw[wire] (r1)--(o);
}


\subsection{Tầng S2: Dịch phải căn chỉnh 31 bit và mặt nạ nhiệt tạo jam}

Tầng S2 (\texttt{posit\_add\_align\_comb}, Hình~\ref{fig:adder-s2}) ghép số hạng
nhỏ thành từ 31 bit $\texttt{y\_seed} = \{1'\mathrm{b0}, 1'\mathrm{b1}, \texttt{y\_frac}, 3'\mathrm{b000}\}$
và dịch phải qua bộ dịch \texttt{dyn\_right\_shifter\_fixed4} theo
\texttt{shift\_amount[4:0]}. Song song với bộ dịch, mạch tổ hợp gồm 25 bộ so sánh
hằng dựng mặt nạ nhiệt $\texttt{discard\_mask}[j] = (j + 3 < \texttt{shift\_amount})$
để thu gọn toàn bộ các bit bị dịch rơi xuống dưới bit~0 thành \texttt{raw\_tail}.

Khi \texttt{both\_nonzero=1} và không bị kẹp, 30 bit cao lấy phần dịch
\texttt{shifted\_full[30:1]}; ngược lại (\texttt{clamped=1} hoặc có một số hạng
bằng 0), 30 bit cao được xóa về 0. Bit \texttt{y[0]} được hợp nhất cùng phần dư
$\texttt{both\_nonzero} \land (\texttt{clamped} \lor \texttt{shifted\_full}[0] \lor \texttt{raw\_tail})$,
còn cờ \texttt{alignment\_tail} lưu riêng trạng thái rơi bit của phép căn chỉnh để
truyền tiếp qua S3 tới S4. Nhờ kiểm soát qua \texttt{both\_nonzero}, nhánh kẹp
không bao giờ tự phát sinh bit jam từ toán hạng 0.

\rtlfig{Datapath chi tiết tầng S2 (\sig{posit\_add\_align\_comb}): dịch phải căn chỉnh 31 bit, mặt nạ nhiệt 25 bit tính song song và hợp nhất \sig{jam}/\sig{alignment\_tail}.}{fig:adder-s2}{
\node[pin] (in1) at (-4.9,.86) {\sig{s1\_y\_frac\_q[25:0]}, \sig{s1\_shift\_q[4:0]}};
\node[pin] (in2) at (.2,.86) {\sig{s1\_clamped\_q}, \sig{s1\_both\_nz\_q}};
\node[pin] (in3) at (5.25,.86) {\sig{s1\_x\_frac\_q[25:0]}, \sig{s1\_sub\_q},\\ \sig{s1\_sf\_q[10:0]}, dấu, cờ};

\node[unit,text width=43mm] (sh) at (-4.9,-1.38) {\hw{Dịch phải 31 bit}\\
 \sig{dyn\_right\_shifter\_fixed4}\\
 $\sig{y\_seed}=\{0,1,f_Y,3'\mathrm{b0}\}$\\
 $\sig{shifted\_full}=\sig{y\_seed}\gg\mathrm{shift}$};
\node[unit,text width=44mm] (msk) at (.2,-1.38) {\hw{Mặt nạ nhiệt 25 bit song song}\\
 $\sig{discard\_mask}[j]=(j+3<\mathrm{shift})$\\
 $\sig{raw\_tail}=\bigvee(f_Y[24:0]\land\sig{mask})$};
\node[unit,text width=36mm] (pas) at (5.25,-1.38) {\hw{Giữ nguyên toán hạng X}\\
 \sig{x\_frac[25:0]}, \sig{subtract},\\
 \sig{sf\_base[10:0]}, \sig{sign\_base},\\
 \sig{both\_zero}, \sig{is\_nar}, \sig{flags\_lo}};

\draw[wire] (in1)--(sh);
\draw[wire] (in2)--(msk);
\draw[wire] (in3)--(pas);
\node[tapdot] at (-4.9,.18) {};
\draw[wire] (-4.9,.18) -- (-1.35,.18) -- (msk.north -| -1.35,0);
\node[tag] at (-2.9,.18) {\sig{f\_Y, shift}};

\node[unit,text width=43mm] (yhi) at (-4.9,-3.85) {\hw{Chốt 30 bit cao của Y}\\
 $\sig{y[30:1]}=(\sig{bnz}\land\neg\sig{clp})$\\
 $\text{ ? }\sig{shifted\_full[30:1]} : 30'\mathrm{d0}$};
\node[unit,text width=45mm] (jam) at (.2,-3.85) {\hw{Tạo jam Y[0] và tail}\\
 $\sig{y[0]}=\sig{bnz}\land(\sig{clp}\lor\sig{sh[0]}\lor\sig{raw\_tail})$\\
 $\sig{tail}=\sig{bnz}\land(\sig{clp}\lor\sig{raw\_tail})$};

\draw[wire] (sh)--(yhi);
\draw[wire] (msk)--(jam);
\coordinate (sh0p) at (-1.45,-2.66);
\node[tapdot] at (-4.9,-2.66) {};
\draw[wire] (-4.9,-2.66) -- (sh0p) -- (jam.north -| sh0p);
\node[tag] at (-3.05,-2.66) {\sig{sh[0]}};

\node[ff,text width=140mm] (r2) at (.15,-5.68) {\hw{S2}\enspace
 \sig{s2\_x\_frac\_q[25:0]} \enspace|\enspace
 \sig{s2\_y\_q[30:0]} \enspace|\enspace
 \sig{s2\_tail\_q}\\
 \sig{s2\_sub\_q}, \sig{s2\_sign\_q}, \sig{s2\_sf\_q[10:0]} \enspace|\enspace
 \sig{s2\_both\_zero\_q}, \sig{s2\_nar\_q}, \sig{s2\_flags\_q[3:0]}};

\draw[wire] (yhi.south)--(r2.north -| yhi);
\draw[wire] (jam.south)--(r2.north -| jam);
\draw[wire] (pas.south)--(r2.north -| pas);

\node[pin] (o) at (.15,-6.82) {Tới S3 (\sig{posit\_add\_arith\_comb}): cộng/trừ 31 bit \sig{X +/- Y} và tính song song \sig{sf\_inc}};
\draw[wire] (r2)--(o);
}


\subsection{Tầng S3 và S4: Cộng/trừ độ lớn, tiền tăng scale và dò bit dẫn đầu}

Hình~\ref{fig:adder-tail} minh họa đầy đủ đường dữ liệu của hai tầng S3
(\texttt{posit\_add\_arith\_comb}) và S4 (\texttt{posit\_add\_lod\_comb}):
\begin{itemize}
  \item \textbf{Tầng S3:} Ghép $X = \{1'\mathrm{b0}, 1'\mathrm{b1}, \texttt{x\_frac}, 3'\mathrm{b000}\}$
  rồi tính \texttt{magnitude[30:0] = subtract ? (X - Y) : (X + Y)}. Vì S1 đã bảo đảm
  $|X|\ge|Y|$ trên mọi cặp hữu hạn khác 0, hiệu $X-Y$ luôn không âm nên không cần
  mạch lấy trị tuyệt đối ở đầu ra S3. Song song với phép cộng/trừ, S3 tính trước
  $\texttt{sf\_inc} = \texttt{sf\_base} + 11'\mathrm{sd1}$ để chuẩn bị cho bước chuẩn hóa ở S5.
  \item \textbf{Tầng S4:} Bộ dò ưu tiên 31 bit \texttt{lod\_lzd\_core\_fixed2} quét từ
  MSB của \texttt{s3\_mag\_q[30:0]} để xác định \texttt{leading\_count[4:0]} và cờ
  \texttt{nonzero}. Đồng thời, nếu phép cộng có tràn nhớ lên bit~30 và bit~0 bằng~1
  ($\texttt{lost\_carry} = M[30]\land M[0]$), S4 jam ngay bit này vào vị trí bit~1 của
  $\texttt{norm\_seed[29:0]} = \{M[29:2], M[1]\lor\texttt{lost\_carry}, M[0]\}$ và OR
  vào \texttt{sticky\_raw} $= \texttt{s3\_tail\_q} \lor \texttt{lost\_carry}$.
\end{itemize}

\rtlfig{Datapath chi tiết tầng S3 (\sig{posit\_add\_arith\_comb}) và tầng S4 (\sig{posit\_add\_lod\_comb}): cộng/trừ độ lớn 31 bit, tiền tăng scale, dò LOD 31 bit và bảo toàn phần dư \sig{lost\_carry}.}{fig:adder-tail}{
\node[pin] (in1) at (-4.95,.76) {\sig{s2\_x\_frac\_q[25:0]}, \sig{s2\_y\_q[30:0]}, \sig{s2\_sub\_q}};
\node[pin] (in2) at (.20,.76) {\sig{s2\_sf\_q[10:0]}};
\node[pin] (in3) at (5.30,.76) {\sig{s2\_tail\_q}, dấu, cờ \sig{Z/N}, \sig{flags\_q}};

\node[unit,text width=41mm] (alu) at (-4.95,-1.05) {\hw{S3: cộng/trừ độ lớn 31 bit}\\
 $X=\{0,1,f_X,3'\mathrm{b0}\}$\\
 $M=\sig{subtract}\text{ ? }(X-Y):(X+Y)$};
\node[unit,text width=44mm] (inc) at (.20,-1.05) {\hw{S3: tiền tăng scale song song}\\
 $\sig{sf\_inc}[10:0]=\sig{sf\_base}+11'\mathrm{sd1}$};
\node[unit,text width=36mm] (m3) at (5.30,-1.05) {\hw{S3: giữ tail \& metadata}\\
 \sig{tail}, \sig{sign}, \sig{both\_zero},\\ \sig{is\_nar}, \sig{flags\_lo[3:0]}};

\draw[wire] (in1)--(alu);
\draw[wire] (in2)--(inc);
\draw[wire] (in3)--(m3);

\node[ff,text width=140mm] (r3) at (.15,-2.72) {\hw{S3}\enspace
 \sig{s3\_mag\_q[30:0]} \enspace|\enspace
 \sig{s3\_sf\_inc\_q[10:0]} \enspace|\enspace
 \sig{s3\_tail\_q}, \sig{s3\_sign\_q}, \sig{s3\_both\_zero\_q}, \sig{s3\_nar\_q}, \sig{s3\_flags\_q[3:0]}};

\draw[wire] (alu.south)--(r3.north -| alu);
\draw[wire] (inc.south)--(r3.north -| inc);
\draw[meta] (m3.south)--(r3.north -| m3);

\node[unit,text width=40mm] (lod) at (-4.95,-5.12) {\hw{S4: bộ dò LOD 31 bit}\\
 \sig{lod\_lzd\_core\_fixed2}\\
 $M[30:0]\to K[4:0],\ \sig{nonzero}$};
\node[unit,text width=46mm] (seed) at (.20,-5.12) {\hw{S4: seed 30 bit \& sticky\_raw}\\
 $\sig{lost\_carry}=M[30]\land M[0]$\\
 $\sig{seed}=\{M[29:2], M[1]{\lor}L, M[0]\}$\\
 $\sig{sticky\_raw}=\sig{tail}\lor\sig{lost\_carry}$};
\node[unit,text width=35mm] (m4) at (5.30,-5.12) {\hw{S4: giữ scale \& cờ}\\
 \sig{sf\_inc[10:0]}, \sig{sign},\\ \sig{both\_zero}, \sig{is\_nar}, \sig{flags\_lo}};

\draw[wire] (r3.south -| lod)--(lod.north);
\draw[wire] (r3.south -| seed)--(seed.north);
\draw[meta] (r3.south -| m4)--(m4.north);

\node[tapdot] at (-4.95,-3.65) {};
\draw[wire] (-4.95,-3.65) -- (-1.45,-3.65) -- (seed.north -| -1.45,0);
\node[tag] at (-3.15,-3.65) {$M[30:0]$};

\node[tapdot] at (5.30,-3.65) {};
\draw[meta] (5.30,-3.65) -- (1.85,-3.65) -- (seed.north -| 1.85,0);
\node[tag] at (3.65,-3.65) {\sig{s3\_tail\_q}};

\node[ff,text width=140mm] (r4) at (.15,-7.05) {\hw{S4}\enspace
 \sig{s4\_count\_q[4:0]}, \sig{s4\_nz\_q} \enspace|\enspace
 \sig{s4\_seed\_q[29:0]}, \sig{s4\_sticky\_q}\\
 \sig{s4\_sf\_inc\_q[10:0]} \enspace|\enspace
 \sig{s4\_sign\_q}, \sig{s4\_both\_zero\_q}, \sig{s4\_nar\_q}, \sig{s4\_flags\_q[3:0]}};

\draw[wire] (lod.south)--(r4.north -| lod);
\draw[wire] (seed.south)--(r4.north -| seed);
\draw[meta] (m4.south)--(r4.north -| m4);

\node[pin] (o) at (.15,-8.18) {Tới S5 (\sig{posit\_add\_norm\_comb}): dịch trái \sig{seed} theo $K$, tính $\sig{sf\_inc}-K$ và phân loại đầu ra};
\draw[wire] (r4)--(o);
}


Một trường hợp dễ kiểm là $\widehat P$ và C có cùng magnitude,
trái dấu và $\Delta sf=0$. Hai significand thẳng hàng,
hiệu bằng 0, LOD không tìm thấy bit 1 và Packer chọn mã Zero.
Nếu số hạng thứ hai chỉ nhỏ hơn một bit thấp, LOD thấy bit 1
ở xa đầu, scale giảm và bit thấp được kéo lên. Vì vậy một
thay đổi nhỏ trước adder có thể làm sai số tương đối tăng
đáng kể trong nhóm gần triệt tiêu. Đó là lý do Chương~\ref{ch:measurements}
dùng sai số tuyệt đối cho nhóm này.

\subsection{Tầng S5: Dịch trái chuẩn hóa, cập nhật scale và phân loại đầu ra}

Khối tổ hợp S5 (\texttt{posit\_add\_norm\_comb}, Hình~\ref{fig:adder-s5}) nhận
bộ đếm \texttt{leading\_count[4:0]} để dịch trái từ \texttt{norm\_seed[29:0]} qua
bộ dịch \texttt{dyn\_left\_shifter\_fixed6}. Vì phép dịch trái 30 bit tự đẩy bit ẩn~1 lên
vượt khỏi biên bit~28 của kết quả trong mọi chế độ (tràn nhớ $K=0$, bình thường
$K=1$, và triệt tiêu sâu $K\ge 2$), một phép trừ duy nhất
$\texttt{sf} = \texttt{sf\_inc} - K$ xử lý thống nhất cả ba chế độ mà không cần bộ
chọn tăng/giảm scale riêng.

Tại thanh ghi đầu ra S5, RTL chỉ chốt 29 bit fraction cao vào \texttt{frac\_hi\_q[28:0]}
và nối thêm 24 bit~0 hằng tại biên giao diện (\texttt{frac = \{frac\_hi\_q, 24'b0\}}),
nhờ vậy loại bỏ được 24 flip-flop luôn mang giá trị 0. Đồng thời, logic phân loại
\texttt{Zero/NaR} đặt tại tầng Packer thay vì chèn 41 MUX xóa bus số học trên đường
găng S5; tầng S5 chỉ cập nhật trực tiếp cờ ngoại lệ (\texttt{flags\_out[4]=is\_nar\_in},
và bit \texttt{inexact} chỉ OR thêm \texttt{sticky\_raw} khi kết quả khác Zero).

\rtlfig{Datapath chi tiết tầng S5 (\sig{posit\_add\_norm\_comb}): dịch trái chuẩn hóa 30 bit, trừ scale một bước và phân loại ngoại lệ không xóa bus dữ liệu.}{fig:adder-s5}{
\node[pin] (in1) at (-4.95,.82) {\sig{s4\_seed\_q[29:0]}, \sig{s4\_count\_q[4:0]}};
\node[pin] (in2) at (0,.82) {\sig{s4\_sf\_inc\_q[10:0]}};
\node[pin] (in3) at (4.95,.82) {\sig{s4\_nz\_q}, \sig{sticky}, dấu, \sig{Z/N}, \sig{flags\_lo}};

\node[unit,text width=41mm] (lsh) at (-4.95,-1.65) {\hw{Dịch trái chuẩn hóa 30 bit}\\
 \sig{dyn\_left\_shifter\_fixed6}\\
 $\{\sig{shifted\_hi[28:0]},\sig{lsb}\}=\sig{seed}\ll K$\\
 $\sig{frac\_hi}=\sig{shifted\_hi}$ (không xóa 0)};
\node[unit,text width=41mm] (sf) at (0,-1.65) {\hw{Cập nhật scale một bước}\\
 $\sig{count\_ext}=\{6'\mathrm{b0},K[4:0]\}$\\
 $\sig{sf[10:0]}=\sig{sf\_inc}-\sig{count\_ext}$\\
 gộp tràn $K{=}0$ và triệt tiêu $K{\ge}2$};
\node[unit,text width=43mm] (cls) at (4.95,-1.65) {\hw{Phân loại Zero/NaR \& cờ}\\
 $\sig{fz}=\neg\sig{nar\_in}\land(\sig{both\_zero}\lor\neg\sig{nz})$\\
 $\sig{sign}=\sig{nar\_in}\text{ ? }1:(\sig{fz}\text{ ? }0:\sig{sign\_base})$\\
 $\sig{inexact}|=\neg\sig{fz}\land\sig{sticky\_raw}$};

\draw[wire] (in1)--(lsh);
\draw[wire] (in2)--(sf);
\draw[wire] (in3)--(cls);
\node[tapdot] at (-4.95,.14) {};
\draw[wire] (-4.95,.14) -- (-1.35,.14) -- (sf.north -| -1.35,0);
\node[tag] at (-2.95,.14) {$K[4:0]$};

\node[ff,text width=142mm] (r5) at (0,-3.92) {\hw{S5}\enspace
 \sig{frac\_hi\_q[28:0]} \enspace|\enspace
 \sig{sf[10:0]} \enspace|\enspace
 \sig{sticky} \enspace|\enspace
 \sig{sign}, \sig{is\_zero}, \sig{is\_nar}, \sig{flags\_out[4:0]}, \sig{out\_valid}};

\draw[wire] (lsh.south)--(r5.north -| lsh);
\draw[wire] (sf.south)--(r5.north -| sf);
\draw[meta] (cls.south)--(r5.north -| cls);

\node[unit,text width=106mm] (pad) at (0,-5.32) {\hw{Khớp giao diện Packer (tiết kiệm 24 Flip-Flop hằng 0)}\\ $\sig{frac[52:0]}=\{\sig{frac\_hi\_q[28:0]},\ 24'\mathrm{b0}\}$};
\draw[wire] (r5.south)--(pad.north);

\node[pin] (o) at (0,-6.50) {Tới Packer P1 (\sig{posit\_pack\_prepare}):\\
 \sig{frac[52:0]}, \sig{sf[10:0]}, \sig{sticky}, \sig{sign}, \sig{is\_zero}, \sig{is\_nar}, \sig{flags\_out}};
\draw[wire] (pad.south)--(o.north);
}


\subsection{Truy vết thanh ghi bộ cộng qua S1--S5}

Để đối chiếu trực tiếp với hai phép tính tay tại Chương~\ref{ch:arithmetic}
($3{,}5 + 0{,}75 = 4{,}25$ và $3{,}5 + (-3{,}25) = 0{,}25$),
Bảng~\ref{tab:adder-trace-examples} ghi lại giá trị chốt trong các thanh ghi
của \texttt{posit\_add} tại ranh giới từng tầng \texttt{S1..S5}. Cùng một bộ
dò LOD 31 bit (\texttt{K=0} khi có tràn nhớ tại bit~30; \texttt{K=4} khi
bit dẫn đầu rơi xuống bit~26 do triệt tiêu) và cùng một phép trừ
\texttt{sf = sf\_inc - K} xử lý thống nhất cả hai chế độ mà không cần bộ
phân nhánh tăng/giảm scale riêng.

\begin{engineeringtable}{Truy vết giá trị thanh ghi qua năm tầng của \texttt{posit\_add} cho trường hợp tràn nhớ ($4{,}25$) và triệt tiêu ($0{,}25$).}{tab:adder-trace-examples}
\begin{tabularx}{\linewidth}{@{}>{\raggedright\arraybackslash}p{16mm} >{\raggedright\arraybackslash}p{62mm} >{\raggedright\arraybackslash}p{62mm}@{}}
\toprule
Tầng & Ca 1: $3{,}5 + 0{,}75 = 4{,}25$ (cộng tràn nhớ) & Ca 2: $3{,}5 + (-3{,}25) = 0{,}25$ (trừ triệt tiêu) \\
\midrule
Đầu vào &
$A{:}\ sf{=}1,\ f{=}\texttt{26'h3000000},\ s{=}0$\newline
$B{:}\ sf{=}-1,\ f{=}\texttt{26'h2000000},\ s{=}0$ &
$A{:}\ sf{=}1,\ f{=}\texttt{26'h3000000},\ s{=}0$\newline
$B{:}\ sf{=}1,\ f{=}\texttt{26'h2800000},\ s{=}1$ \\
\addlinespace[3pt]
\texttt{S1} &
$\sig{choose\_a}{=}1$, $\sig{s1\_shift\_q}{=}2$, $\sig{s1\_sub\_q}{=}0$, $\sig{s1\_sf\_q}{=}1$ &
$\sig{choose\_a}{=}1$, $\sig{s1\_shift\_q}{=}0$, $\sig{s1\_sub\_q}{=}1$, $\sig{s1\_sf\_q}{=}1$ \\
\texttt{S2} &
$X{=}\texttt{31'h38000000}$, $\sig{s2\_y\_q}{=}\texttt{31'h0c000000}$, $\sig{tail}{=}0$ &
$X{=}\texttt{31'h38000000}$, $\sig{s2\_y\_q}{=}\texttt{31'h34000000}$, $\sig{tail}{=}0$ \\
\texttt{S3} &
$\sig{s3\_mag\_q}{=}\texttt{31'h44000000}$ (bit~30 $=1$), $\sig{s3\_sf\_inc\_q}{=}2$ &
$\sig{s3\_mag\_q}{=}\texttt{31'h04000000}$ (bit~26 $=1$), $\sig{s3\_sf\_inc\_q}{=}2$ \\
\texttt{S4} &
$\sig{s4\_count\_q}{=}0$, $\sig{s4\_seed\_q}{=}\texttt{30'h04000000}$, $\sig{sticky}{=}0$ &
$\sig{s4\_count\_q}{=}4$, $\sig{s4\_seed\_q}{=}\texttt{30'h04000000}$, $\sig{sticky}{=}0$ \\
\texttt{S5} &
$\sig{sf}{=}2-0=2$, $\sig{frac\_hi\_q}{=}\texttt{29'h02000000}$ ($1{,}0625\times 2^2=4{,}25$) &
$\sig{sf}{=}2-4=-2$, $\sig{frac\_hi\_q}{=}\texttt{29'h00000000}$ ($1{,}0\times 2^{-2}=0{,}25$) \\
\bottomrule
\end{tabularx}
\end{engineeringtable}

\subsection{Packer P1 tạo G/R/S và phân loại}

Quá trình đóng gói hai tầng của bộ cộng/tích được mô tả trong
Algorithm~\ref{alg:posit-pack}, tương ứng với hai mô-đun
\texttt{posit\_pack\_prepare} (P1) và \texttt{posit\_pack\_finish} (P2).

\begin{reportpseudocode}{alg:posit-pack}{Đóng gói và làm tròn RNE hai tầng (Two-Stage Posit(32,3) RNE Packing)}
  {Kết quả chuẩn hóa $(s, sf[10:0], f[52:0], \mathrm{sticky}, z, \eta, \mathrm{flags_{in}}[4:0])$.}
  {Từ mã đóng gói Posit(32,3) $d[31:0]$ và cờ trạng thái $\mathrm{flags}[4:0]$.}
  \begin{tabularx}{\linewidth}{@{}r >{\raggedright\arraybackslash}X >{\raggedright\arraybackslash}p{.33\linewidth}@{}}
    \algline{1}{\textbf{P1} \quad $r_0\gets \neg sf[10]$, $\mathrm{offset}[4:0]\gets sf[7:3]\oplus\{5\{sf[10]\}\}$}{Xác định bit regime đầu và độ dịch không rẽ nhánh.}
    \algline{2}{$\mathrm{payload}[58:0]\gets \{r_0, \neg r_0, sf[2:0], f[52:0], 1'b0\}$}{Ghép 59 bit gồm regime, bit dừng, mũ và phân số.}
    \algline{3}{$\mathrm{shifted}\gets (\mathrm{payload}\gg_r \mathrm{offset})$}{Dịch phải điền $r_0$ qua 5 tầng MUX để tạo regime.}
    \algline{4}{$m_{\mathrm{trunc}}[30:0]\gets \mathrm{shifted}[58:28]$, $G\gets \mathrm{shifted}[27]$, $R\gets \mathrm{shifted}[26]$}{Tách độ lớn 31 bit và hai bit làm tròn Guard, Round.}
    \algline{5}{$S\gets \mathrm{sticky}\lor (\bigvee f[24:0])\lor \bigvee_{k=0}^{30}(\{sf[2:0],f[52:25]\}_k \land (k<\mathrm{offset}))$}{Tính bit Sticky song song bằng mặt nạ nhiệt kế.}
    \algline{6}{$\mathrm{sel}\gets \mathrm{Classify}(\eta, z, sf, f, \mathrm{sticky})\in\{\mathrm{NAR},\mathrm{ZERO},\mathrm{MAX},\mathrm{MIN},\mathrm{NORM}\}$}{Ưu tiên NaR, Zero và bão hòa tại biên $sf\in[-240,240]$.}
    \algline{7}{\textbf{P2} \quad $\mathrm{round\_up}\gets G\land (m_{\mathrm{trunc}}[0]\lor R\lor S)$}{Kiểm tra điều kiện làm tròn gần nhất chẵn (RNE).}
    \algline{8}{$\mathrm{base}\gets \{1'b0, m_{\mathrm{trunc}}\}\oplus \{32\{s\}\}$, $\mathrm{inc}\gets s\oplus \mathrm{round\_up}$}{Gộp bù hai khi $s=1$ và lượng tăng RNE thành một bộ cộng.}
    \algline{9}{$d_{\mathrm{round}}\gets \mathrm{base}+\mathrm{inc}$}{Tính từ mã đã làm tròn trên nhánh bình thường.}
    \algline{10}{\textbf{if} $\mathrm{sel}=\mathrm{NAR}$ \textbf{then} $d\gets \texttt{32'h80000000}$, $\mathrm{flags}\gets \texttt{5'b10000}$}{Xuất mẫu NaR và đặt cờ ngoại lệ NaR.}
    \algline{11}{\textbf{else if} $\mathrm{sel}\in\{\mathrm{ZERO},\mathrm{MAX},\mathrm{MIN}\}$ \textbf{then} $d\gets \mathrm{Const}(\mathrm{sel},s)$}{Xuất $0$, $\pm\mathrm{maxpos}$ hoặc $\pm\mathrm{minpos}$ theo $\mathrm{sel}$.}
    \algline{12}{\textbf{else} $d\gets \mathrm{Clamp}(m_{\mathrm{trunc}}, s, d_{\mathrm{round}})$, $\mathrm{flags}[1]\gets \mathrm{flags_{in}}[1]\lor G\lor R\lor S$}{Chặn tràn biên max/min và cập nhật cờ \texttt{inexact}.}
    \algline{13}{\textbf{return} $d, \mathrm{flags}$}{Chốt kết quả Posit(32,3) cuối đường ống P2.}
  \end{tabularx}
\end{reportpseudocode}

P1 nhận scale signed 11 bit, fraction 53 bit, dấu, sticky và
metadata ngoại lệ. Bit regime đầu là nghịch đảo dấu của scale.
Với $k\ge0$, offset bằng phần thấp của $k$. Với $k<0$, offset
là $|k|-1$ trên miền normal. Payload 59 bit ghép bit regime đầu,
terminator, exponent 3 bit, fraction 53 bit và một bit 0.
Bộ dịch phải điền bit regime đầu vào các vị trí trống bên trái,
tạo magnitude giữ lại từ \texttt{shifted[58:28]}. Hai bit
tiếp theo là guard và round.

\rtlfig{Packer cuối: P1 dựng magnitude và G/R/S song song với phân loại; P2 ghép RNE với bù hai và chọn mã. Mọi đường đều đi qua hai tầng valid/ready.}{fig:packer-stages}{
\node[pin] (i) at (0,.69) {\sig{sf[10:0]}, \sig{frac[52:0]}, sticky, dấu, Zero/NaR, flags 5};
\node[unit,text width=42mm] (asm) at (-5,-1.46) {\hw{Regime + payload}\\$r=\neg sf[10]$\\$o=sf[7:3]\oplus\{5\{sf[10]\}\}$\\$P=\{r,\neg r,sf[2:0],f,0\}$};
\node[unit,text width=42mm] (mask) at (0,-1.46) {\hw{Mặt nạ sticky 31 bit}\\$U=\{sf[2:0],f[52:25]\}$\\chọn $U_j$ nếu $j<o$\\$S=\mathrm{sticky}\lor |f[24:0]\lor |\mathrm{mask}$};
\node[unit,text width=42mm] (class) at (5,-1.46) {\hw{Phân loại song song}\\NaR $\to$ Zero $\to$ MAX $\to$ MIN\\MAX: $sf\ge240$; MIN: $sf<-240$\\\sig{special\_sel}, flags};
\foreach \n in {asm,mask,class}{\draw[wire] (i.south)--++(0,-.24)-|(\n.north);}
\node[unit,text width=42mm] (sh) at (-5,-3.91) {\hw{Dịch phải 59 bit}\\$W=P\gg o$; điền $r$\\$m=W[58:28]$\\$G=W[27]$, $R=W[26]$};
\draw[wire] (asm)--node[tag,right]{59, $o$:5, $r$:1}(sh);
\draw[wire] (asm.east)--(mask.west);
\node[ff,text width=145mm] (r1) at (0,-5.6) {\hw{P1}\enspace $m[30:0]$\enspace|\enspace $G,R,S$\enspace|\enspace dấu\enspace|\enspace \sig{special\_sel[2:0]}, flags};
\draw[wire] (sh)--(r1.north -| sh);
\draw[wire] (mask)--node[tag,right]{$S$:1}(r1.north -| mask);
\draw[meta] (class)--node[tag,right]{loại, flags}(r1.north -| class);
\node[unit,text width=66mm] (r) at (-3.73,-7.34) {\hw{RNE + mã dấu: một phép cộng}\\
 $u=G\land(m[0]\lor R\lor S)$\\
 $D_{\mathrm{r}}=(\{0,m\}\oplus\{32\{s\}\})+(s\oplus u)$};
\node[unit,text width=62mm] (end) at (3.73,-7.34) {\hw{Biên và cờ}\\
 test $m=0$ / $m=\mathtt{7fff\,ffff}$\\
 chọn min/max theo dấu\\nhánh thường: inexact OR G/R/S};
\draw[wire] (r1.south -| r)--(r);
\draw[wire] (r1.south -| end)--(end);
\node[pick,text width=81mm] (mux) at (0,-9.38) {\hw{Chọn mã theo}~\sig{special\_sel}~\hw{và biên}};
\draw[wire] (r.south)--++(0,-.32)-|(mux.north -| r.east);
\draw[wire] (end.south)--++(0,-.32)-|(mux.north -| end.west);
\node[ff,text width=95mm] (r2) at (0,-10.65) {\hw{P2} \quad \sig{d[31:0]}, \sig{flags[4:0]}, \sig{out\_valid}};
\draw[wire] (mux)--(r2);
}



P1 tính sticky song song với bộ dịch. Fraction thấp
\texttt{frac[24:0]} luôn nằm dưới round, nên OR trực tiếp với
sticky đầu vào. Các bit cao hơn chỉ tham gia sticky khi vị
trí của chúng nhỏ hơn offset tương ứng. Mỗi bit dùng một so
sánh hằng theo offset và cổng AND với dữ liệu. Cuối cùng OR
các bit được chọn. Nhờ vậy phép tìm sticky không đi nối tiếp
qua toàn bộ đường barrel shifter.

\subsection{Packer P2 làm tròn, kẹp và mã hóa dấu}

Với magnitude giữ $m$ và quyết định tăng $r\in\{0,1\}$,
đầu ra dương bình thường là $m+r$. Với số âm, biểu thức
$-(m+r)=\mathop{\mathrm{NOT}}m+(1-r)$ trên 32 bit cho phép
ghép việc tăng RNE và mã hóa bù hai. RTL chọn XOR theo dấu
và dùng increment \texttt{sign XOR round\_up}. Đường chọn
ngoại lệ vẫn có thể trả Zero, NaR, minpos hoặc maxpos.

P1 mang theo chỉ dấu nhánh đặc biệt và flags. NaR ưu tiên,
sau đó Zero và các nhánh biên. Với nhánh thường,
\texttt{inexact} OR thông tin mất từ trước với G/R/S.
\texttt{approx\_cut} vẫn chỉ báo còn số hạng tại điểm dừng
bộ nhân. Việc tách hai cờ giúp scoreboard phân biệt
lượng tử hóa ở Packer với giới hạn vòng.

\section{Phạm vi đối chứng NF và quire}

Các khối P1/P2, nhân lặp, adder và Packer cuối của MAC-NF và MAC-F được
triển khai đồng bộ trên cùng bộ nguồn RTL dùng chung khối nhân. NF có đường tích đi qua
Packer tích hai tầng P1/P2 (\texttt{posit\_product\_pack}) và Parser tích
(\texttt{posit\_magnitude\_parse}) rồi đến adder (Hình~\ref{fig:nf-f-split}).
Cụ thể, ở tầng P1, mạch tổ hợp \texttt{posit\_pack\_prepare} dịch phải 59 bit
và tạo mặt nạ nhiệt sticky rồi chốt \texttt{p1\_mag\_trunc\_q[30:0]},
\texttt{G/R/S}, \texttt{p1\_sign\_q}, \texttt{p1\_special\_sel\_q[2:0]},
\texttt{p1\_flags\_q[4:0]} và \texttt{p1\_valid\_q}. Ở tầng P2, mạch tổ hợp
\texttt{posit\_pack\_finish} làm tròn RNE 32 bit (\texttt{sign=0}) và ghép nhánh
ngoại lệ rồi chốt \texttt{d\_q[31:0]}, \texttt{flags\_q[4:0]}, \texttt{sign\_q}
và \texttt{p2\_valid\_q}. Tiếp đó, \texttt{posit\_magnitude\_parse} giải mã tổ hợp
trước khi bank tích NF chốt \texttt{norm\_valid\_q} để đưa vào S1.

\rtlfig{Biên tích trước S1 của MAC-F và MAC-NF. Hai cột biểu diễn hai thiết kế riêng; cả hai giữ C từ lúc Parser retire, còn \sig{posit\_product\_pack} gồm hai tầng pipeline P1/P2 rõ mạch tổ hợp và FF.}{fig:nf-f-split}{
\node[unit,text width=116mm] (src) at (0,0) {\hw{Adapter sau nhân:}~\sig{mul\_prepack\_norm\_comb}\\
 \sig{acc[25]} chọn $(sf+1,\{\sig{acc[24:0]},0\})$ hoặc $(sf,\{\sig{acc[23:0]},00\})$};
\node[head] (hf) at (-3.5,-1.18) {MAC-F};
\node[head] (hn) at (3.65,-1.18) {MAC-NF};
\node[unit,text width=66mm] (p1c) at (3.65,-2.52) {\hw{P1 tổ hợp:}~\sig{posit\_pack\_prepare}\\
 dịch phải 59 bit (\sig{frac}=\{26, 27'b0\}, \sig{sticky}=0)\\mặt nạ nhiệt sticky; phân loại \sig{special\_sel}};
\node[ff,text width=66mm] (p1r) at (3.65,-4.48) {\hw{P1 FF:}~\sig{posit\_product\_pack} (\sig{p1\_ready})\\
 \sig{p1\_mag\_trunc\_q[30:0]}, \sig{G/R/S}, \sig{p1\_sign\_q}\\\sig{p1\_special\_sel\_q[2:0]}, \sig{p1\_flags\_q[4:0]}, \sig{p1\_valid\_q}};
\node[unit,text width=66mm] (p2c) at (3.65,-6.42) {\hw{P2 tổ hợp:}~\sig{posit\_pack\_finish} (\sig{sign=0})\\
 \sig{round\_up} RNE 32 bit, MUX \sig{special\_sel}/biên, cờ \sig{inexact}};
\node[ff,text width=66mm] (p2r) at (3.65,-8.18) {\hw{P2 FF:}~\sig{posit\_product\_pack} (\sig{p2\_ready})\\
 \sig{d\_q[31:0]} (magnitude), \sig{flags\_q[4:0]}, \sig{sign\_q}, \sig{p2\_valid\_q}};
\node[unit,text width=66mm] (mp) at (3.65,-10.02) {\sig{posit\_magnitude\_parse}\\XOR + sentinel $\to$ LOD 31\\tái tạo frac 26, $sf$ 11, $s,Z,N$};
\node[ff,text width=66mm] (nr) at (3.65,-11.96) {\hw{Bank tích NF} (\sig{product\_capture})\\\sig{norm\_valid\_q}, \sig{product\_frac\_q}, \sig{norm\_sf\_q}, $s/Z/N$, \sig{norm\_flags\_q}};
\node[unit,text width=54mm] (af) at (-3.5,-13.62) {\hw{S1 nhận tích F và C}};
\node[unit,text width=66mm] (an) at (3.65,-13.62) {\hw{S1 nhận tích NF và C}};
\draw[wire] (src.south) -- ++(0,-.26) -| (hn);
\draw[wire] (src.south) -- ++(0,-.26) -| (hf);
\draw[wire] (hn)--node[tag,right]{\sig{core\_done}}(p1c);
\draw[wire] (p1c)--(p1r);
\draw[wire] (p1r)--(p2c);
\draw[wire] (p2c)--(p2r);
\draw[wire] (p2r)--node[tag,right]{magnitude 32 + \sig{sign\_q}}(mp);
\draw[wire] (mp)--node[tag,right]{\sig{product\_capture}}(nr);
\draw[wire] (hf)--(af) node[pos=.28,tag,right]{frac 26 và $sf$ 11\\đưa thẳng vào S1\\với \sig{core\_done}};
\draw[wire] (nr)--(an);
\draw[meta] (p2r.west)--++(-.68,0) |- (nr.west) node[pos=.42,tag,left]{\sig{flags\_q[4:0]}};
}
 Khi giải thích
chênh lệch tài nguyên, so hai kết quả trong cùng một lượt
tổng hợp. Khi giải thích chu kỳ, dùng cùng dãy A/B và cùng
điều kiện backpressure; không lấy độ sâu hình khối để suy
ra tần số hay throughput.

Với quy ước đo tuần tự từng phép một (mỗi lần 1 phép đang chạy, \texttt{out\_ready=1}, đếm gộp từ cạnh chấp nhận đầu vào đến khi hoàn tất rút kết quả và \texttt{in\_ready} sẵn sàng ở cạnh kế tiếp), số chu kỳ của một phép MAC khác 0 phụ thuộc số bit~1 được chọn $t = \mathrm{popcount}(X_{\mathrm{mant}})$ ở bộ nhân như sau:
\begin{itemize}
\item \textbf{MAC-F:} $T_F(t) = 13 + \max(1,t) \in [14, 25]$ chu kỳ với $0 \le t \le 12$ (Parser \texttt{P1/P2}: 2; nạp ngữ cảnh/\texttt{T1}: 3; lặp \texttt{T2/T3}: $\max(1,t)$; Adder \texttt{S1..S5}: 5; Packer \texttt{P1/P2}: 2; chuyển tiếp retire sang \texttt{in\_ready}: 1; riêng khoảng cách từ cạnh nhận đến \texttt{out\_valid} là $12 + \max(1,t)$ cạnh).
\item \textbf{MAC-NF:} $T_{NF}(t) = 16 + \max(1,t) = T_F(t) + 3 \in [17, 28]$ chu kỳ. Do \texttt{posit\_magnitude\_parse} là mạch tổ hợp, phần chênh đúng $+3$ chu kỳ đến từ 2 tầng FF \texttt{P1/P2} của \texttt{posit\_product\_pack} và 1 tầng FF \texttt{norm\_valid\_q} tại bank tích NF trước tầng cộng.
\item \textbf{MAC-Quire:} Lệnh \texttt{CLEAR} mất $3$ chu kỳ; lệnh \texttt{MAC} với tích khác 0/NaR mất $36 + \max(1,t) = T_F(t) + 23 \in [37, 48]$ chu kỳ do FSM đi qua 30 chu kỳ xử lý tại bank tích lũy 27 từ sau \texttt{core\_done} (thay vì 7 chu kỳ nối thẳng qua Adder/Packer ở MAC-F), còn tích bằng 0 hoặc NaR bỏ qua vòng bank 27 bước để xong sau 1 chu kỳ; lệnh \texttt{READ} mất $59$ chu kỳ để quét/giảm 27 từ quire và đóng gói \texttt{posit\_pack}.
\end{itemize}

Quire dùng mô-đun đỉnh \texttt{posit\_mac\_quire} và các module backend
cập nhật theo từng từ 36 bit. Lệnh MAC căn tích
sang lưới $2^{-480}$ và giữ $Q$ có dấu 972 bit. CLEAR đưa
$Q$/count/status về trạng thái khởi tạo; READ quét $Q$ để
tạo fraction và sticky cho Packer. Chi tiết hiệu quả tối ưu hóa của kiến trúc
backend theo word (từ phiên bản dịch toàn vector sang cập nhật theo từ 36 bit)
và số chu kỳ đo trên chuỗi tích lũy được trình bày cụ thể trong Chương~\ref{ch:measurements}.

\rtlfig{Backend quire cập nhật theo từ 36 bit. MAC cập nhật từng word rồi commit Q; READ đọc Q đã commit, dò bit dẫn đầu theo word và làm tròn đúng một lần.}{fig:quire-context}{
\node[head] at (-3.85,0) {Ghi: lệnh MAC};
\node[head] at (3.85,0) {Đọc: lệnh READ};
\node[ff,text width=63mm] (prod) at (-3.85,-1.15) {\hw{Chốt tích từ core}\\magnitude 26 Q24, $sf$ 11, dấu/cờ};
\node[unit,text width=63mm] (al) at (-3.85,-3.24) {\hw{Căn lưới $2^{-480}$}\\
 $h=sf+456$; hệ số signed 27 bit\\
  index $=h_+\ \mathrm{div}\ 36$; offset $=h_+\bmod36$\\
  $h_+=\max(h,0)$ trên nhánh hợp lệ};
\draw[wire] (prod)--(al);
\node[unit,text width=63mm] (w) at (-3.85,-5.53) {\hw{Chọn term của word đang xử lý}\\
 dịch hệ số trong 63 bit; word 36 bit\\
 chọn 0, lát thấp, lát cao hoặc dấu};
\draw[wire] (al)--node[tag,right]{hệ số 27, index 5, offset 6}(w);
\node[unit,text width=63mm] (sum) at (-3.85,-7.55) {\hw{Cộng theo word 0 $\to$ 26}\\
 \sig{q\_words[step]} + term + carry\\
 tổng 37 bit; bit 36 truyền vòng sau};
\draw[wire] (w)--(sum);
\node[ff,text width=63mm] (pend) at (-3.85,-9.43) {\hw{Bank pending}\\
 ghi word tổng; commit khi xong word 26 và không overflow};
\draw[wire] (sum)--node[tag,right]{tổng thấp 36}(pend);
\node[ff,text width=64mm] (q) at (3.85,-1.15) {\hw{Q đã commit}\\27 word $\times$ 36 bit; signed 972 bit};
\draw[wire] (pend.south)--(-3.85,-10.26)--(.28,-10.26) -- (.28,-1.15) -- (q.west);
\node[tag,rotate=90] at (.28,-9.15) {commit};
\draw[wire] ($(q.west)+(0,.32)$)--(-.24,-.83)|-(sum.east);
\node[unit,text width=63mm] (fetch) at (3.85,-3.02) {\hw{FETCH: lấy độ lớn word}\\
 chọn word theo step từ thấp lên cao\\
 XOR theo dấu Q, cộng carry $\to$ word 36 bit};
\draw[wire] (q)--(fetch);
\node[ff,text width=63mm] (abs) at (3.85,-4.78) {\hw{Chốt}~\sig{abs\_word\_q}~\hw{và carry}};
\draw[wire] (fetch)--(abs);
\node[unit,text width=63mm] (red) at (3.85,-6.54) {\hw{REDUCE: dò và tạo cửa sổ}\\
 ưu tiên bit 1 cao nhất trong word\\
 cửa sổ 3 word dịch $35-\mathrm{lead\_bit}$};
\draw[wire] (abs)--(red);
\draw[enable] (red.east)--(7.56,-6.54) |- (fetch.east);
\node[ff,text width=63mm] (out) at (3.85,-8.67) {\hw{Cập nhật khi word $\ne0$}\\
 $sf=36\,\mathrm{step}+\mathrm{lead\_bit}-480$\\
 frac 53; sticky OR phần thấp\\và các word cũ};
\draw[wire] (red)--(out);
\node[unit,text width=63mm] (p) at (3.85,-10.52) {\hw{Packer RNE sau word 26}};
\draw[wire] (out)--(p);
\node[pin] (r) at (3.85,-11.6) {\sig{d[31:0]}, flags, \sig{out\_cmd=READ}};
\draw[wire] (p)--(r);
\node[explain,text width=63mm] at (-3.85,-11.6) {CLEAR xóa Q và trạng thái.\\CLEAR/MAC trả ack với $d=0$.};
}



\subsection{Commit quire và phép đọc giữ trạng thái}

Một MAC hợp lệ tạo term trên lưới quire từ dấu, magnitude
Q24 và scale. Phần căn tích chuẩn bị descriptor gồm hệ số,
chỉ số word và offset. Backend cập nhật theo từng word với
carry/borrow qua ranh giới, rồi commit state hoàn chỉnh.
Trong thời gian cập nhật, giao diện giữ một lệnh outstanding.
MAC chỉ trả acknowledge sau khi state đã cập nhật và count
tăng. Nếu sink kéo ready thấp, response giữ ổn định và không
commit lặp. C được giữ ở driver đối với NF/F; quire có state
trong backend và không có cổng C.

READ xác định dấu và bit dẫn đầu của Q, dựng một cửa sổ
fraction và tổng hợp sticky cho các bit thấp hơn cửa sổ.
Packer dùng dữ liệu đó để làm tròn đầu ra. READ giữ Q,
count và trạng thái loss của chuỗi. CLEAR tạo epoch mới
với Q và count bằng 0. Khi count đã đạt 1.024, lệnh MAC
tiếp bị chặn; READ và CLEAR vẫn có vai trò kết thúc hoặc
khởi tạo lại chuỗi. Điểm này được sử dụng trong benchmark:
mỗi chuỗi dài 1.024 có đúng 1.024 MAC rồi READ.

\begin{engineeringtable}{Giao diện lệnh của \texttt{posit\_mac\_quire} theo đặc tả kỹ thuật.}{tab:quire-command-flow}
\begin{tabularx}{\linewidth}{@{}>{\raggedright\arraybackslash}p{22mm}Y Y@{}}
\toprule
Lệnh & Điều kiện và tác động state & Response \\
\midrule
CLEAR & Xóa Q, count và trạng thái của epoch. & Ack với out\_cmd=CLEAR, d=0. \\
MAC & Nhận A/B và cấu hình khi còn dung lượng, tạo và commit term. & Ack với d=0, flags phản ánh trạng thái/loss hiện tại. \\
READ & Quét state Q đã commit, giữ state. & D là kết quả RNE và flags của READ. \\
\bottomrule
\end{tabularx}
\end{engineeringtable}
